% Lösungen Einheit 5
\einheitenkopf{Einheit 5 von 5 · Kurvenuntersuchung im Sachzusammenhang}
\erg{39}{a) A \quad b) B \quad c) C \quad d) D \quad e) B \quad f) A}
\erg{40}{a) nimmt zu, immer langsamer \quad b) nimmt ab, immer schneller \quad c) nimmt ab, immer langsamer \quad d) nimmt zu, immer schneller \quad e) bei $t = 4$, also um 10 Uhr (Wendepunkt) \quad f) $W(4) = 6$, also 600 m³}
\erg{41}{a) $h'(t) = -1{,}5t^2 + 6t = -1{,}5t\,(t - 4) = 0 \Leftrightarrow t = 0$ oder $t = 4$; $h(4) = 16$, Ränder $h(0) = h(6) = 0$: nach 4 s ist die Drohne mit 16 m am höchsten \quad b) $N'(t) = -0{,}3t^2 + 3t$, $N'(10) = 0$, $N''(t) = -0{,}6t + 3$, $N''(10) = -3 < 0$: Hochpunkt; $N(10) = 70$, Ränder $N(0) = 20$, $N(14) = 39{,}6$: größte Nutzerzahl 70\,000 nach 10 Monaten \quad c) $K'(t) = 50\,(1 - 0{,}5t)\,\mathrm{e}^{-0{,}5t}$, $K''(t) = 50\,(0{,}25t - 1)\,\mathrm{e}^{-0{,}5t} = 0 \Leftrightarrow t = 4$ ($K''$ wechselt von $-$ nach $+$); $K'(4) = -50\,\mathrm{e}^{-2} \approx -6{,}77$: nach 4 Stunden nimmt die Konzentration am stärksten ab, um etwa 6,77 mg/l pro Stunde \quad d) $f'(t) = 1{,}5(t - 2)(t - 6)$; $t = 6$ liegt außerhalb und wird verworfen; $f(2) = 36$, $f(0) = 20$, $f(5) = 22{,}5$: höchster Wasserstand 36 cm (nach 2 Stunden), niedrigster 20 cm (zu Beginn, am Rand)}
\erg{42}{a) Nullstellen $0$ und $6$: Breite $6 \cdot 5$ cm $= 30$ cm $\le 32$ cm; $s'(x) = 0{,}3x\,(x - 4)$, Tiefpunkt $(4 \mid -3{,}2)$: Tiefe $3{,}2 \cdot 5$ cm $= 16$ cm $> 15$ cm – der Trog passt nicht \quad b) $r'(t) = -0{,}6t^2 + 6t$, $r'(5) = 15$: Nach 5 Minuten wächst die Zuflussrate am schnellsten, um 15 Liter pro Minute je Minute. \quad c) (1) $d$ ist der Höhenunterschied von A gegenüber B; (2) Zeitpunkte, an denen sich der Unterschied gerade nicht ändert; (3) bei $t = 10$ ist der Unterschied lokal am größten; (4) er beträgt 10 cm und ist größer als an den Rändern (0 cm und 7,2 cm). Aufgabenstellung: „Ermitteln Sie, wann Pflanze A der Pflanze B am weitesten voraus ist, und geben Sie den Vorsprung an.“ – nach 10 Wochen, 10 cm.}
\erg{43}{a) Paul setzt $T' = 0$; das liefert die Stelle der höchsten Temperatur (14 Uhr), nicht der stärksten Zunahme. Richtig: $T''(t) = -0{,}6t + 1{,}8 = 0 \Leftrightarrow t = 3$; $T'(3) = 2{,}7$: um 11 Uhr steigt die Temperatur am stärksten, um 2,7 °C pro Stunde \quad b) $Z''(t) = -1{,}2t + 3{,}6 = 0 \Leftrightarrow t = 3$; $Z'(3) = 5{,}4$: um 17 Uhr nimmt die Besucherzahl am stärksten zu, um 540 Besucher pro Stunde}
\erg{44}{a) $V'(t)$ ist der Grenzwert der Quotienten $\frac{\Delta V}{\Delta t}$: Liter geteilt durch Minuten. \quad b) Die Temperatur steigt bis zum Ende der Messung; der größte Wert liegt dann am Rand des Zeitraums, und dort muss die Tangente nicht waagerecht sein.}
\erg{45}{a) Nullstellen $\pm 2$: Breite 4 m, Höhe $z(0) = 2$ m, Länge 3 m; $V = 4 \cdot 2 \cdot 3 = 24$ m³ \quad b) $r'(x) = 0{,}3(x - 2)(x - 4)$; $r(2) = 3$, $r(4) = 2{,}6$, Ränder $r(0) = 1$, $r(6) = 4{,}6$: größter Radius 4,6 cm am Rand; Block $6 \text{ cm} \times 9{,}2 \text{ cm} \times 9{,}2 \text{ cm}$, $V = 507{,}84$ cm³, Masse $\approx 355{,}5$ g}
\erg{46}{a) $h''(x) = -0{,}00012x + 0{,}006 = 0 \Leftrightarrow x = 50$; $h'(50) = 0{,}15$ (an den Rändern $h' = 0$): am steilsten nach 50 m mit 15\,\% \quad b) $\tan\alpha = 0{,}15$, $\alpha \approx 8{,}5^\circ \le 10^\circ$: die Vorgabe ist eingehalten}
\erg{47}{a) zwischen 200 und 800 Stück ($2 < x < 8$) \quad b) Fixkosten von 8000 €, die auch ohne Produktion anfallen \quad c) 3 m}
